\chapter*{Lời Mở Đầu}
\addcontentsline{toc}{chapter}{Lời Mở Đầu}
\markboth{Lời Mở Đầu}{Lời Mở Đầu}

\begin{truyen}{Lời Mở Đầu}{Opening Words}
\chuhoa{C}{ó thời con người mở điện thoại khi cần.}
\textit{(There was a time when people opened the phone when necessary.)}

Bây giờ rất nhiều người sống như thể mình chỉ tạm đặt đời thật xuống để quay lại với màn hình cho kịp.
\textit{(Now many people live as if real life is only something briefly put down so they can hurry back to the screen.)}

Học sinh học cùng điện thoại.
\textit{(Students study with the phone beside them.)}

Phụ huynh ăn cùng điện thoại.
\textit{(Parents eat with the phone beside them.)}

Người lớn đi bộ, tập thể dục, lái xe, vào nhà vệ sinh, thậm chí ngồi cạnh nhau mà mắt vẫn bị hút xuống một cái màn hình nhỏ.
\textit{(Adults walk, exercise, drive, go to the toilet, and even sit beside one another while their eyes are still pulled down by a small screen.)}

Cuốn này không viết để chửi điện thoại một cách ngu ngốc.
\textit{(This book is not written to curse the phone in a simplistic way.)}

Nó viết để soi cho rõ hơn: điện thoại có ích thật, nhưng ở rất nhiều người, phần ích chỉ chiếm một góc nhỏ còn phần hại đã lan ra gần hết đời sống.
\textit{(It is written to make one thing clear: the phone has real usefulness, but in many people the useful part occupies only a small corner while the harmful part has spread across almost all of life.)}

Muốn tỉnh ngộ thì không được phiến diện.
\textit{(If we want awakening, we must not be one-sided.)}

Nhưng cũng không được hèn đến mức lấy vài ưu điểm nhỏ để tha bổng cho cả một nếp sống đang bị nuốt mất.
\textit{(But we must not be so weak that we pardon an entire devoured way of life just because of a few small advantages.)}
\end{truyen}

\begin{baihoc}
Điện thoại không chỉ lấy thời gian.
\textit{(A phone does not only take time.)}

Nó lấy từng mẩu chú ý rất nhỏ.
\textit{(It takes tiny pieces of attention.)}

Và chính những mẩu rất nhỏ đó, khi mất quá lâu, sẽ làm con người không còn chịu được đời sống không có kích thích ngay lập tức.
\textit{(And those tiny pieces, when lost for too long, make a person unable to endure a life without immediate stimulation.)}
\end{baihoc}

\begin{tuhocnhanh}
1. Một ngày của em đang bị cắt vụn bởi điện thoại ở những khoảnh khắc nào? \textit{(At which moments is your day being sliced into pieces by the phone?)}

2. Em đang dùng điện thoại như công cụ, hay đang để nó dùng sự chú ý của em như nguyên liệu? \textit{(Are you using the phone as a tool, or letting it use your attention as raw material?)}

3. Nếu phải nói thật, em sợ điều gì nhất ở những lúc không có màn hình bên cạnh? \textit{(If you are honest, what do you fear most in moments when no screen is beside you?)}
\end{tuhocnhanh}

\ngancach